\documentclass[../main]{subfiles}
\begin{document}
\chapter{Cavitación}
\img{images/07/cavitacion.jpg}{Cavitación en alábes de bomba}

\info{
Contenido}{
Cavitación
}

\section{Cavitación}
La cavitación se produce cuando la presión antes de entrar en el rotor es menor a la presión de vapor del fluido. Esto provoca la evaporación del mismo en forma de pequeñas burbujas, las cuales se condensan rápida y violentamente unos instantes después cuando la presión aumenta drásticamente.

\img{images/07/cavitacion2.jpg}{Fenómeno de cavitación}

Cuando en una tubería que transporta agua la presión interior comienza a caer por debajo de la presión atmosférica se inicia el vacío. A medida que desciende la presión, el vacío aumenta, aunque el agua sigue líquida. Es a partir de un determinado valor de vacío cuando el agua cambia de estado y se transforma en vapor.

Un líquido puede pasar a vapor debido a dos circunstancias, o bien porque aumente su temperatura o bien porque disminuya su presión. En el caso del agua, y para una temperatura de 10ºC, considerada como la temperatura media del agua en los bombeos, el cambio de estado se produce cuando la presión absoluta es de 9 mm de columna de mercurio, lo que equivale a:

Cuando el agua pasa a vapor, libera gases que se encontraban disueltos y se forman unas burbujas (las cuales se denominan cavidades, de ahí el nombre del fenómeno) que viajan con la corriente líquida y que, al ganar presión nuevamente aguas abajo, explotan con violencia. Es lo que se conoce como cavitación.

Si, por ejemplo, una válvula en una conducción opera con una escasa abertura y un alto caudal provoca una alta pérdida de presión debido al rozamiento de la corriente líquida con el obturador de la válvula. Al mismo tiempo se produce un notable incremento de la velocidad del fluido, consecuencia directa de la disminución del área de paso ya que:


La hidrodinámica es la rama de la física que estudia el comportamiento de los fluidos en movimiento. Es una parte crucial de la mecánica de fluidos y tiene aplicaciones importantes en ingeniería, física y otras disciplinas científicas. Entre los fenómenos que estudia la hidrodinámica se encuentra la cavitación, un proceso que puede causar daños significativos en sistemas hidráulicos y propulsores marinos.

La cavitación se produce cuando la presión en un fluido desciende por debajo de su presión de vapor, provocando la formación de burbujas de vapor. Estas burbujas, al colapsar, generan ondas de choque y fuerzas intensas que pueden dañar superficies metálicas y componentes de bombas, turbinas y hélices.

\subsection{Fórmulas Fundamentales}

A continuación, se presentan algunas de las fórmulas fundamentales de la hidrodinámica que serán utilizadas en este trabajo práctico:

\info{Ecuación de Bernoulli}{
\begin{equation}
P + \frac{1}{2} \rho v^2 + \rho gh = \text{constante}
\end{equation}
Donde:
\begin{itemize}
    \item \( P \) es la presión del fluido (Pa)
    \item \( \rho \) es la densidad del fluido (kg/m³)
    \item \( v \) es la velocidad del fluido (m/s)
    \item \( g \) es la aceleración debido a la gravedad (9.81 m/s²)
    \item \( h \) es la altura relativa en el campo gravitacional (m)
\end{itemize}
}

\info{Número de Reynolds}{
\begin{equation}
Re = \frac{\rho v L}{\mu}
\end{equation}
Donde:
\begin{itemize}
    \item \( Re \) es el número de Reynolds (adimensional)
    \item \( \rho \) es la densidad del fluido (kg/m³)
    \item \( v \) es la velocidad del fluido (m/s)
    \item \( L \) es una longitud característica (m)
    \item \( \mu \) es la viscosidad dinámica del fluido (Pa·s)
\end{itemize}
}

\info{Ecuación de Continuidad}{
\begin{equation}
A_1 v_1 = A_2 v_2
\end{equation}
Donde:
\begin{itemize}
    \item \( A_1 \) y \( A_2 \) son las áreas de las secciones transversales 1 y 2 (m²)
    \item \( v_1 \) y \( v_2 \) son las velocidades del fluido en las secciones 1 y 2 (m/s)
\end{itemize}
}

La cavitación, como se mencionó anteriormente, ocurre cuando la presión en un punto del fluido cae por debajo de la presión de vapor del fluido. La ecuación de Bernoulli es fundamental para entender cómo las variaciones de presión pueden llevar a la cavitación.

\section{Consecuencias de la Cavitación}

La cavitación puede causar serios problemas en sistemas hidráulicos. Los impactos de las burbujas colapsantes pueden erosionar metales y otros materiales, reduciendo la eficiencia y vida útil de componentes como hélices, turbinas y bombas. Las consecuencias más comunes de la cavitación incluyen:

\begin{itemize}
    \item Daños en superficies metálicas: Las ondas de choque y los microchorros producidos por el colapso de burbujas pueden erosionar y picar superficies metálicas.
    \item Pérdida de eficiencia: La cavitación puede causar fluctuaciones de presión y flujo que afectan el rendimiento de bombas y turbinas.
    \item Ruido y vibraciones: El colapso de burbujas produce ruido y vibraciones que pueden ser perjudiciales para la integridad de los equipos.
\end{itemize}

\section{Prevención de la Cavitación}

Para prevenir la cavitación, es importante mantener la presión del fluido por encima de su presión de vapor en todas las partes del sistema. Algunas estrategias comunes incluyen:

\begin{itemize}
    \item Diseñar sistemas con geometrías que minimicen la caída de presión.
    \item Aumentar la presión de entrada del fluido.
    \item Utilizar materiales resistentes a la erosión por cavitación.
    \item Controlar la velocidad del fluido para evitar condiciones que favorezcan la cavitación.
\end{itemize}

\section{Cuestionario de Investigación}

\begin{enumerate}
    \item ¿Qué es la hidrodinámica y cuáles son sus principales aplicaciones?
    \item ¿Cómo se define la cavitación y qué condiciones la provocan?
    \item ¿Cuáles son las consecuencias más comunes de la cavitación en sistemas hidráulicos?
    \item ¿Qué estrategias se pueden implementar para prevenir la cavitación en bombas y turbinas?
    \item ¿Cómo afecta la presión de vapor del fluido a la cavitación?
    \item ¿Qué papel juega la ecuación de Bernoulli en el estudio de la cavitación?
    \item ¿Cuál es la relación entre el número de Reynolds y la cavitación?
    \item ¿Cómo se puede utilizar la ecuación de continuidad para analizar sistemas hidráulicos?
    \item ¿Qué tipos de materiales son más resistentes a la erosión por cavitación?
    \item ¿Cómo se puede medir y monitorear la cavitación en un sistema hidráulico?

\end{enumerate}

\section{Problemas a Resolver}

\begin{enumerate}
    \item En una tubería de 0.05 m² de sección transversal, el fluido se mueve a 2 m/s. Si la sección transversal se reduce a 0.02 m², ¿cuál será la nueva velocidad del fluido?
    \item Una bomba eleva agua a una altura de 10 m con una presión de entrada de 200 kPa. ¿Cuál es la presión del agua en la salida de la bomba, suponiendo que no hay pérdidas de energía?
    \item Un fluido con una densidad de 1000 kg/m³ fluye a través de una tubería con un diámetro de 0.1 m a una velocidad de 3 m/s. ¿Cuál es el número de Reynolds si la viscosidad del fluido es 0.001 Pa·s?
    \item Si en un sistema hidráulico la velocidad del fluido es 4 m/s en una sección de 0.03 m², ¿cuál será la velocidad en una sección de 0.01 m²?
    \item La presión de vapor del agua a 20°C es 2.3 kPa. Si la presión en una zona del sistema desciende a 1.5 kPa, ¿se producirá cavitación?
    \item Una hélice gira en agua con una velocidad de 5 m/s y la presión atmosférica es 101 kPa. Si la presión en las palas de la hélice baja a 95 kPa, ¿existe riesgo de cavitación?
    \item En una tubería, el área de entrada es 0.04 m² y la velocidad del fluido es 1 m/s. Si el área de salida es 0.02 m², calcula la velocidad del fluido en la salida.
    \item Un fluido tiene una densidad de 850 kg/m³ y se mueve a una velocidad de 2 m/s en una tubería de 0.05 m de diámetro. ¿Cuál es el número de Reynolds si la viscosidad del fluido es 0.002 Pa·s?
    \item Calcula la presión a una altura de 5 m en una columna de agua, si la presión en la base es 250 kPa.
    \item Si la presión de un fluido en una sección de una tubería es 150 kPa y la velocidad es 3 m/s, calcula la presión en una sección donde la velocidad aumenta a 5 m/s, considerando una densidad de 1000 kg/m³ y despreciando cambios en la altura.

\end{enumerate}

\end{document}
